% -*- root: ../thesis.tex -*-
%!TEX root = ../thesis.tex
% ******************************* Thesis Chapter 5 ****************************


% ----------------------- paths to graphics ------------------------

% change according to folder and file names
\graphicspath{{chapter_drosophila/figures/}}
% ----------------------- Drosophila chapter ------------------------

% Could also remove section heading here
\section{Context within thesis}

This chapter shows how the consolidation of odor-behavior associations through individual synaptic pathways in the mushroom body of Drosophila can be seen as an instantiation of the parallel pathway theory (PPT). Interestingly though, the Hebbian-like learning in these pathways is non-local to the pre- and postsynaptic partners of a synapse, mediated via recurrently connected Dopamine neurons that gate plasticity. This non-locality allows to rephrase the PPT in the mushroom body as a recurrent neural network acting on memory representations evolving over time scales of memory consolidation.

%\section{Contribution}
%The work presented in this chapter was done based on collaboration with Desire Laber, Tania Fernandez del Valle Alquicira, Lisa Scheunemann and David Owald, who we discussed experimental data with and who gave feedback on the models. I developed the model, performed the simulations and analyzed the data. This work has not been published.
%
%\section{Introduction}
%
%\section{Results}
%
%\section{Methods}
%
%\section{Discussion}

\section{Chapter summary}
This chapter illustrated how the parallel pathway theory (PPT) can be implemented in the mushroom body of Drosophila and how it can be re-interpreted as a recurrent neural network acting on memory representations evolving over timescales of memory consolidation. This perspective will be the basis for the next chapter, where we will generalize it as a phenomenological model for brain-wide engram dynamics under systems consolidation.



%\begin{figure}[htb]
%  \centering
%
%\resizebox{1\textwidth}{!}{\input{5/pgf/Exemplary_ROC.pgf}}
%
%  \caption[Example of a ROC curve]{ Example of using python to generate a pgf-figure which has the same fonts as the main latex document. Run python plot\_exemplary\_roc.py from the Python directory to generate the pgf-file. }
%  \label{fig:ROC_curve}
%\end{figure}  




% ---------------------------------------------------------------------------
%: ----------------------- end of thesis sub-document ------------------------
% ---------------------------------------------------------------------------

